\part{Classical Fields and Geometry}

\chapter{Maxwell fields are differential forms}
\label{ch:forms}

\physicalquestion{How can electric charge and magnetic flux be described in
a way that every observer can integrate consistently?}

\modelassumptions{Spacetime is a smooth oriented four-manifold.  The
electromagnetic field is classical, sources are smooth on the region under
study, and singular charges or material interfaces are excluded unless they
are represented by removed subsets or distributions.}

\section{What can be pulled back and integrated?}

A smooth $n$-manifold is a Hausdorff, second-countable space locally modeled
on $\R^n$ by smoothly compatible charts.  Its tangent and cotangent bundles
are $TM$ and $T^*M$.  A differential $p$-form is a smooth section of
$\Lambda^pT^*M$; write $\Omega^p(M)$ for their space.

The wedge product is alternating and the exterior derivative
$d:\Omega^p(M)\to\Omega^{p+1}(M)$ is the unique graded derivation extending
the differential of functions and satisfying $d^2=0$.  For a smooth map
$f:N\to M$, pullback obeys $f^*d=df^*$.  Integration pairs compactly
supported top-degree forms with oriented manifolds.  Stokes' theorem states
\[
\int_Xd\omega=\int_{\partial X}\omega.
\]

Closed forms lie in $\ker d$, exact forms in $\im d$, and
$H^p_{\mathrm{dR}}(M)=\ker d/\im d$ records the obstruction to finding a
global potential.

\section{Which conservation law is structural?}

\begin{theorem}[Maxwell charge conservation]
Let $F\in\Omega^2(M)$ be the electromagnetic field and let
$J\in\Omega^3(M)$ be the current.  If a metric-dependent Hodge star satisfies
the Maxwell equations
\[
dF=0,\qquad d{\star F}=J,
\]
then $dJ=0$.  For every compact oriented four-region $X$,
\[
\int_{\partial X}J=0.
\]
\end{theorem}

\begin{proof}
Apply $d$ to the sourced equation: $dJ=d^2{\star F}=0$.  Stokes' theorem then
gives $\int_{\partial X}J=\int_XdJ=0$.
\end{proof}

\section{How does this return to charge measurements?}

The boundary of a spacetime slab consists of two spatial slices and a timelike
side.  The theorem says that the difference of total charge on the slices
equals the current flux through the side.  This conclusion does not depend on
a coordinate continuity equation; that equation is a chart expression of
$dJ=0$.  The physical model still assumes that $J$ includes every relevant
source.  Apparent nonconservation may instead signal current crossing an
unmodeled boundary or failure of the classical description.

\section{What further contact should be proved?}

\begin{exercise}[Flux depends only on homology]
If $F$ is closed and two oriented $2$-cycles differ by the boundary of a
$3$-chain, prove that their integrals of $F$ agree.  State the dual dependence
on the de Rham class $[F]$.
\end{exercise}

\begin{exercise}[Gauge potentials are local data]
If $F=dA$ on a contractible region, prove that replacing $A$ by $A+d\lambda$
does not change $F$.  Give an example on a punctured region where a closed
one-form is not exact.
\end{exercise}


\chapter{Metrics turn geometry into dynamics}
\label{ch:metric}

\physicalquestion{What additional structure turns invariant field forms into
lengths, wave equations, freely falling paths, and tidal acceleration?}

\modelassumptions{The manifold carries a smooth nondegenerate metric of fixed
signature.  Test bodies are small enough to neglect their backreaction and,
when geodesic motion is invoked, non-gravitational forces are absent.}

\section{What does a metric determine without coordinates?}

A pseudo-Riemannian metric is a smooth nondegenerate symmetric bilinear form
$g$ on each tangent space.  An orientation and $g$ determine a volume form
$\operatorname{vol}_g$ and a Hodge star by
\[
\alpha\wedge\star\beta
=\langle\alpha,\beta\rangle_g\operatorname{vol}_g.
\]

A connection on $TM$ assigns covariant derivatives $\nabla_XY$, is linear in
$X$, satisfies a Leibniz rule in $Y$, and permits vectors at nearby points to
be compared infinitesimally.  It is torsion free when
$\nabla_XY-\nabla_YX=[X,Y]$ and metric compatible when $\nabla g=0$.

\section{Why is there a canonical metric connection?}

\begin{theorem}[Levi--Civita]
Every pseudo-Riemannian metric has a unique torsion-free, metric-compatible
connection.  It is determined by the coordinate-free Koszul formula
\begin{align*}
2g(\nabla_XY,Z)={}&Xg(Y,Z)+Yg(Z,X)-Zg(X,Y)\\
&-g(X,[Y,Z])+g(Y,[Z,X])+g(Z,[X,Y]).
\end{align*}
\end{theorem}

\begin{proof}
Metric compatibility expands each derivative of an inner product into two
covariant-derivative terms.  Add the expansions for $Xg(Y,Z)$ and
$Yg(Z,X)$, subtract the one for $Zg(X,Y)$, and use zero torsion to replace
differences of covariant derivatives by Lie brackets.  This yields the
formula and hence uniqueness by nondegeneracy of $g$.  Conversely the formula
defines $\nabla_XY$ uniquely; direct substitution verifies linearity, the
Leibniz rule, zero torsion, and metric compatibility.
\end{proof}

\section{How do connection and star return to physics?}

The Hodge star in Chapter~\ref{ch:forms} supplies the metric part of Maxwell's
equations.  The Levi--Civita connection defines geodesics by
$\nabla_{\dot\gamma}\dot\gamma=0$ and curvature by
\[
R(X,Y)Z=\nabla_X\nabla_YZ-\nabla_Y\nabla_XZ-\nabla_{[X,Y]}Z.
\]
For a variation of geodesics with separation field $J$, the equation
$\nabla_T\nabla_TJ=R(T,J)T$ turns curvature into relative acceleration.  A
gravitational-wave interferometer uses this tidal effect; it does not measure
coordinate components of the metric in isolation.

\section{What further contact should be proved?}

\begin{exercise}[The Hodge star is unique]
Prove pointwise existence and uniqueness of $\star$ from the induced inner
product on exterior powers.  Determine $\star^2$ in terms of form degree and
metric signature.
\end{exercise}

\begin{exercise}[Killing fields give conserved quantities]
If $X$ satisfies $\mathcal L_Xg=0$ and $\gamma$ is an affinely parametrized
geodesic, prove that $g(X,\dot\gamma)$ is constant.
\end{exercise}


\chapter{Hamiltonian evolution is symplectic}
\label{ch:symplectic}

\physicalquestion{Why does conservative time evolution preserve phase-space
volume even when it stretches and folds individual regions?}

\modelassumptions{The classical system has a finite-dimensional smooth phase
space, its dynamics is deterministic, and the Hamiltonian has no explicit
time dependence.  Dissipation and stochastic forcing are excluded.}

\section{What replaces a metric on phase space?}

A symplectic form is a closed nondegenerate two-form $\omega$ on a manifold
$M$.  Nondegeneracy identifies every one-form with a unique vector field.  A
Hamiltonian $H\in C^\infty(M)$ therefore determines $X_H$ by
\[
\iota_{X_H}\omega=dH.
\]
The Poisson bracket is $\{f,h\}=\omega(X_f,X_h)$.  In canonical coordinates
the definition reproduces Hamilton's equations, but neither $\omega$ nor
$X_H$ depends on those coordinates.

\section{Why can Hamiltonian flow not compress phase volume?}

\begin{theorem}[Liouville]
The local flow $\varphi_t$ of a Hamiltonian vector field preserves the
symplectic form and the Liouville volume $\omega^n/n!$ on a $2n$-dimensional
phase space.
\end{theorem}

\begin{proof}
Cartan's formula and $d\omega=0$ give
\[
\mathcal L_{X_H}\omega
=d(\iota_{X_H}\omega)+\iota_{X_H}d\omega=d^2H=0.
\]
Thus
$\frac d{dt}\varphi_t^*\omega
=\varphi_t^*\mathcal L_{X_H}\omega=0$,
so $\varphi_t^*\omega=\omega$.  Pullback respects wedge products, hence it
also fixes $\omega^n/n!$.
\end{proof}

\section{What physical consequence follows?}

An ensemble transported by a Hamiltonian flow cannot converge into a smaller
phase-volume attractor.  This explains why ordinary friction cannot be
represented by an autonomous Hamiltonian on the same finite-dimensional
phase space: friction contracts volumes and violates the theorem's
assumptions.  In equilibrium statistical mechanics the preserved volume is
the invariant measure from which microcanonical counting begins.

\section{What further contact should be proved?}

\begin{exercise}[Magnetic curvature gives the Lorentz force]
On $T^*Q$ replace the canonical symplectic form by
$\omega_B=\omega_0+q\pi^*B$ for a closed two-form $B$ on $Q$.  For kinetic
Hamiltonian $H$, derive the Lorentz-force equation.
\end{exercise}

\begin{exercise}[Moment maps recover angular momentum]
For the lifted rotation action on $T^*\R^3$, prove that the moment map is
$J(q,p)=q\times p$.  Show directly that a rotation-invariant Hamiltonian
conserves $J$, and compare with Chapter~\ref{ch:lie}.
\end{exercise}


\chapter{Local symmetry becomes a connection}
\label{ch:gauge}

\physicalquestion{How can internal states at different spacetime points be
compared when no preferred frame exists globally?}

\modelassumptions{Internal states form fibers of a smooth vector bundle.  A
compact Lie group $G$ acts on the fibers, local gauge transformations are
changes of frame, and classical gauge fields are smooth connections.}

\section{Which object compares neighboring fibers?}

A vector bundle $E\to M$ is locally a product $U\times V$ with linear
transition maps.  A connection is a linear map
\[
\nabla:\Omega^0(M,E)\to\Omega^1(M,E),
\qquad \nabla(fs)=df\otimes s+f\nabla s.
\]
Extending it to $E$-valued forms defines curvature
$F_\nabla=\nabla^2\in\Omega^2(M,\End E)$.

For a principal $G$-bundle, a local frame represents the connection by a
Lie-algebra-valued one-form $A$, with
\[
F_A=dA+A\wedge A.
\]
A gauge transformation changes $A$ but conjugates $F_A$; curvature is the
local obstruction to path-independent transport.

\section{Which identity constrains every gauge field?}

\begin{theorem}[Bianchi identity]
Every connection satisfies $d_A F_A=0$, or intrinsically
$\nabla F_\nabla=0$.
\end{theorem}

\begin{proof}
In a local frame, $d_A\eta=d\eta+A\wedge\eta-(-1)^p\eta\wedge A$ for a
$p$-form.  Substitute $F_A=dA+A\wedge A$, use $d^2=0$ and the graded Leibniz
rule, and observe that all remaining terms cancel associatively.  Since the
identity transforms covariantly, it is independent of the chosen frame.
\end{proof}

\section{How does local symmetry produce field equations?}

With an invariant inner product on the Lie algebra and the Hodge star from
Chapter~\ref{ch:metric}, the Yang--Mills action is
\[
S(A)=\frac12\int_M\langle F_A\wedge\star F_A\rangle.
\]
A compactly supported variation gives $d_A\star F_A=0$ in vacuum.  The Bianchi
identity is kinematic; the Yang--Mills equation is dynamical.  For $G=U(1)$
they reduce to the two Maxwell equations.  Nonabelian groups additionally
self-couple through $A\wedge A$.

The Standard Model uses a particular gauge group and matter representations.
Those choices are empirical inputs to the model; gauge geometry organizes
their charges and couplings but does not derive the observed representation
list from differential geometry alone.

\section{What further contact should be proved?}

\begin{exercise}[Gauge variation of curvature]
For $A^g=g^{-1}Ag+g^{-1}dg$, prove
$F_{A^g}=g^{-1}F_Ag$ and show that the Yang--Mills action is gauge invariant.
\end{exercise}

\begin{exercise}[Part III synthesis]
On a principal $U(1)$-bundle, connect the four descriptions
\[
\text{potential}\longrightarrow\text{connection}\longrightarrow
\text{curvature}\longrightarrow\text{Maxwell field}.
\]
Derive both Maxwell equations, distinguish the Bianchi identity from the
Euler--Lagrange equation, and identify exactly where orientation, metric, and
bundle topology enter.
\end{exercise}

\parttransition{Curvature and differential equations are local.  They cannot
detect every effect of a flat connection around a hole, nor explain why a
transport coefficient remains an integer under continuous deformation.
Global observables force holonomy, characteristic classes, K-theory, and
index.}
